\setcounter{section}{-1}
\section{Plan}\label{sec:plan}

\begin{enumerate}[leftmargin=2.2em]
\item Section~\ref{sec:overview}: intro, plan, notations
  \begin{itemize}
  \item introduce general framework
  \item random 2 SAT
  \item high level overview of our proof
  \end{itemize}

\item Section~\ref{sec:graph-analysis}: basic results 
  \begin{itemize}
  \item warm up: number of expected solutions
  \item Minimal graph with positive UNSAT probability
  \item Construction of the graph with density $m= 1.5\cdot n$ edges and $\Prob[SAT]=1$.
  \item State and proof the theorem that if $m>1.5\cdot n$, then $\Prob[UNSAT] >0$.
  \item Observation that if graph $G$ has many ($k$) small (constant size at most $C$) barbells, then the probability $\Prob[SAT]=e^{\Omega(-k/C)}$.
  \item Discussion that graphs with many edge-disjoint barbells are easy. Then the next natural question is how to deal with hard graphs. What does it mean that a graph does not have many small barbells?
  Mention that there are graphs $3$ or $4$ regular with large girth $g(G)=\Omega(\log(n))$ that do not have even one small barbell. Such graphs are hard.
  \item New subsection: Regularization of hard graphs. The purpose here is to get some useful regularity properties in case of hard graphs. Specifically, we apply a series of reduction arguments that only increase $\Prob[SAT]$ and keep positive excess and sufficiently large graph size. 
  \begin{itemize}
  	\item More convenient description of easy graphs (those with many small edge-disjoint barbells): the graph has many small edge-disjoint cycles. State the graph-theoretic LEMMA that graph connected graph with $k$ 
	edge disjoint cycles (of any length) has at least $\lceil(k-1)/3\rceil$ edge-disjoint barbells.
	\item COROLLARY for the $\Prob[SAT]=e^{-\Omega(n)}$, when there are linearly many small cycles in total and all connected components have positive excess.
  	\item Reduce largest degree $\Delta(G)$ to a constant $\Delta_0\ge\Delta(G)$ by splitting large degree vertices. This reduction has the property that $\Prob[SAT]$ and the number of vertices $n$ may only increase, 	
	but excess would decrease, however we ensure that it is still positive and linear in $n$ by choosing big enough $\Delta_0(\eps)$.
  	\item Elimination of vertices of degrees $0$ and $1$. Obviously, it only may increase the excess and $\Prob[SAT]$ (probability does not change). 
	It can only reduce the number of vertices by a constant factor. Make it an observation without any special environment (like lemma or claim)
	\item Elimination of connected components with non positive excess. Again the excess and $\Prob[SAT]$ may only increase. Need only to check that we delete at most a constant fraction of vertices. Make it an 
	observation. 
	\item Make one formal CLAIM summarizing the reduction of a hard graph so far: Each vertex has degree $2\le deg(v)\le\Delta_0$; every connected component has positive excess; the overall excess is positive and 
	linear in $n$.
	\item Make another formal CLAIM stating: if $G$ has at least $\rho_1\cdot n$ vertices whose neighborhoods are not trees or at least $\rho_2\cdot n$ loops, then $G$ has $\rho$ edge-disjoint short cycles.
	The proof is easy and short: in a constant degree graph short cycles at each of the vertices with non-tree neighborhood may only intersect (by edges) at most a constant $C(\Delta_0)$ of other such cycles. 
	\item Conclusion. We may assume that the ``hard'' graph has no loops, has degree of each vertex $2\le deg(v)\le\Delta_0$, is locally tree-like with at most $\rho\cdot n$ vertices.
   \end{itemize}     
  \end{itemize}

\item Section~\ref{sec:solution-components}: Expected number of components in
      the solution graph
  \begin{itemize}
  \item Proof of the main identity
  \end{itemize}

\item Section~\ref{sec:difficult-graphs}: the analysis for difficult graphs.
\end{enumerate}
